\chapter{Open Problems and Computational Projects}\label{chap:open}

Each chapter of this book ends at a precise question. This chapter collects
them. We separate three kinds: problems about the Collatz map itself, problems
about Switching MM\index{Switching MM}, and computational projects that are well posed, finite and
likely to repay effort. Where a problem is stated in an earlier chapter we
restate it in its sharpest form and say what is already known. Several problems
have been sharpened by results proved in this book, and we flag where a
natural formulation turns out to be the wrong one.

\section{The divergence half}\label{sec:ch11:divergence}

\begin{problem*}[O1. Ces\`aro equidistribution]\designation{O1}\label{prob:ch11:o1}
Prove \cref{conj:cesaro}: along every Syracuse orbit, $\liminf_N\bar k_N>\log_23$. Bounded orbits satisfy this automatically (\cref{sec:ch9:conj}), so the content is in the unbounded ones. Informally, prove that along no orbit can the running mean of the exponents $\vt(3y_j+1)$ have upper limit at most $\log_23=1+\beta$, which by \cref{cor:ch9:freq} is what an unbounded orbit would need.
\end{problem*}

This is the divergence half of the conjecture itself, restated
(\cref{sec:ch9:conj}). \Cref{prop:cesaro} supplies only the converse implication, which is an
identity, and \cref{rem:ch9:circular} explains why numerical checks of
$\bar k_N>\log_23$ on terminating orbits carry no information. A weaker target
that would still be new is a statement of the form: along every orbit, the
proportion of strides with $k_j=1$ has upper density strictly less than one.
Even this is open, and \cref{cor:ch4:multistride} shows that arbitrarily long
runs of $k_j=1$ do occur.

\begin{problem*}[O2. Trajectory mixing]\designation{O2}\label{prob:ch11:o2}
Fix $a$ and a constant $C>1/|\log(3/4)|$, which is about $3.48$, and let $x_0$ be drawn uniformly from the odd integers up to $X$.
Show that the empirical distribution of $y_j\bmod3^a$ over the strides before the orbit reaches $1$, stopped after $\lfloor C\log X\rfloor$ strides if the orbit has not reached $1$ by then, converges in probability, as $X\to\infty$, to the stationary law $\pi_a$ of \cref{lem:ch9:mixing}.
In the model the cap is reached with probability tending to zero; its only role is to make the empirical distribution defined for every starting value.
\end{problem*}

The problem is posed for random starting values because no version can hold
for every orbit: the orbit of $(4^m-1)/3$ reaches $1$ after one stride and
stays there. As
\cref{lem:ch9:mixing} shows, the model chain on residues modulo $3^a$, in which
the halving exponents are independent Geometric$(\tfrac12)$ draws, mixes\index{mixing}
\emph{exactly} in $a$ strides, with every nontrivial eigenvalue zero, so there
is no spectral\index{spectral gap} question left
for the model. The entire content is the transfer to the map, whose halving
exponents are produced by the orbit itself. Tao\index{Tao, T.} proves a fine scale form of
the transfer for the law of a single iterate of a random starting
value~\cite[Propositions~1.9 and~1.14]{tao2022}; the problem asks for the
empirical distribution along the whole orbit.

The evidence is strong. Along orbits of up to about $25000$ strides from random starting values near $2^{10000}$, the residues modulo $3^5$ approach $\pi_5$ at the rate of independent sampling, and they scatter about that benchmark about as much as the model chain does (\cref{sec:ch9:single}).

\begin{problem*}[O3. Certificates with growing horizon]\designation{O3}\label{prob:ch11:o3}
\Cref{thm:nolyapunov} and \cref{cor:ch4:multistride} rule out certificates\index{certificate} $W(S^{(J)}(x))\le\rho\,W(x)$ with $W$ increasing, for any fixed $J$, and a decreasing $W$ fails at $x=5$ for every horizon, since $S^{(J)}(5)=1$; without monotonicity a
certificate bounded away from zero exists exactly when the conjecture holds
(\cref{rem:ch4:exists,rem:ch4:monotone}). Find an increasing certificate whose
horizon $J(x)$ grows with $x$, say $J(x)=O(\log x)$, with
$W(S^{(J(x))}(x))<W(x)$ for every odd $x\ge3$.
\end{problem*}

For increasing $W$ the condition reads $S^{(J(x))}(x)<x$, so such a certificate is equivalent to the conjecture together with a bound on
stopping times\index{stopping time}: the problem is the conjecture in a stronger form, not a separate question. Its value
is that it separates what is known on average from what is not. The density
one version is already known: by the density theorem\index{density theorem} (\cref{thm:density}),
for every $C>0$ the odd $x$ with $S^{(J)}(x)<x$ for some $J\le C\log x$ have
density one. The corresponding statement for the total stopping time, that
almost every $x$ reaches $1$ within $C\log x$ strides, is open for every $C$;
it would show that almost every integer reaches $1$, and in the model it holds
exactly when $C>1/|\log(3/4)|\approx3.48$.

\begin{problem*}[O3$'$. Potentials of unbounded depth]\designation{O3$'$}\label{prob:ch11:o3p}
\Cref{chap:mmcollatz} shows that every potential $\log x+\varphi(x)$ with
$\varphi$ depending on boundedly many binary digits of $x$ fails, and that the
run length potentials $\log x+c\,\vt(x+1)$ fail at a fraction at least $0.143$
of odd integers. Determine the infimum, over potentials of depth $h$ (reading the runs of the first $h$ blocks), of the fraction of odd integers at which the potential fails to decrease, whether it tends to zero as $h\to\infty$, and if so how fast.
The failure set of such a potential is, up to a set of density zero, a union of residue classes modulo powers of $2$, so its density is the total measure of those classes, which by \cref{thm:terras} is the probability of failure in the Bernoulli model: the problem is a question about the model, and its answer cannot reveal arithmetic structure beyond it.
\end{problem*}

\begin{problem*}[O4. The third cumulant of total stopping times]\designation{O4}\label{prob:ch11:o4}
The census of \cref{sec:ch3:missing} locates the variance deficit: the deficit of about
ten percent is an additive, apparently bounded term of about $2.6$ in $\log^2$
units, concentrated in the deterministic final descent of the orbit, which
indicates that the model is asymptotically right about the variance. The skewness is not explained:
it is $0.67$ at $n\approx2^{29}$ against $0.96$ for the inverse Gaussian law of
the model. Prove that the variance deficit is bounded, and explain the third
cumulant.
As stated, the first part contains the conjecture, since the deficit of a dyadic block $[2^k,2^{k+1})$ is finite only if every integer in it reaches $1$; the question is to prove that the deficit is bounded assuming the conjecture, or unconditionally for the integers of each block whose orbits reach $1$ within a fixed multiple of $k$ steps.
\end{problem*}

\begin{problem*}[O5. Uniform integrability]\designation{O5}\label{prob:ch11:o5}
Settled by \cref{prop:uimartingale}: the iterate martingale\index{martingale} of the accelerated model,
stopped when it first drops to $1$ or below, is not uniformly integrable\index{uniform integrability}, and \cref{chap:path} explains why its deterministic counterpart, the indicator of reaching $1$, leaves nothing to certify short of the conjecture itself. There is no separate open problem here.
\end{problem*}

\section{The cycle half}\label{sec:ch11:cycles}

\begin{problem*}[C1. Arithmetic sharpening of the cycle identity]\designation{C1}\label{prob:ch11:c1}
Find a function $f$, derived from the arithmetic of the parity vector\index{parity vector} rather than the size of the elements, such that every nontrivial cycle with $m$ odd steps, $k$ steps in all and least element $x_{\min}$ satisfies $x_{\min}<m/(3\Lambda f(m))$, where $\Lambda=k\log2-m\log3$ is the value at that cycle; \cref{lem:ch8:window} gives this with $f=1$. By \cref{cor:ch8:window}, $f(m)\ge\frac m3(1+\eta)$ for some fixed $\eta>0$ and all large $m$ would confine the least element of every nontrivial cycle with many odd steps below a constant of about $1/(2\eta)$ for small $\eta$, and so reduce the cycle half to a finite verification; any $f$ bounded below by a constant $b>1$ would lower every threshold of \cref{tab:ch8:thresholds} by at least the factor $b$.
\end{problem*}

An effective irrationality measure\index{irrationality measure} for
$\log_23$ cannot do this on its own (\cref{prop:ch8:insufficient}):
a measure with exponent $\nu$ and constant $c$ excludes only $r<(3cX\log2)^{1/\nu}$. On a cycle
$x(2^k-3^m)=c_k(x)$, where $c_k(x)$ is a sum of terms $3^j2^i$ indexed by the
parity word, and the window lemma\index{window lemma} shows that the least element is confined to
an interval with ends a factor of more than $m/3$ apart.
In the computations of \cref{chap:cycles} the primitive balanced words with $100\le m\le2000$ sit near the top of the window, at about $0.72$ of its upper end (an observation, not a theorem), so arguments that use only the size of the elements, and therefore apply equally to rational cycles, are unlikely to gain the factor there; the problem
is really whether $2^k-3^m$ can divide $c_k(w)$ for a nontrivial word, perhaps
attacked by a block decomposition in the style of the thinning argument that
settled Erd\H{o}s $125$.

\begin{problem*}[C2. Effective measures and the finite range]\designation{C2}\label{prob:ch11:c2}
Find explicit $\nu$ and $c$ with $|\log_23-s/r|\ge c\,r^{-\nu}$ for all integers $s$ and $r\ge1$, with $\nu$ as close to $2$ as possible; the explicit bound of Rhin\index{Rhin, G.}~\cite{rhin1987}, in the form quoted in \cref{sec:ch8:limits}, corresponds to $\nu\approx14.3$. By \cref{rem:irrmeasure} such a pair excludes the cycle lengths $r<(3cX\log2)^{1/\nu}$ without any enumeration, but this range can never exceed the exact ladder value $r_*(X)$, and \cref{tab:ch8:insufficient} shows that even the idealized pair $\nu=2$, $c=0.01773$ falls short of it by factors from about $2$ to more than $50$. The interest of the problem is therefore Diophantine, not the cycle bound.
\end{problem*}

\begin{problem*}[C3. Sharper accountings at the next rungs]\designation{C3}\label{prob:ch11:c3}
Hercher\index{Hercher, C.}'s record $r\ge1.375\times10^{11}$~\cite[Corollary~29]{hercher2023},
unconditional since Ba\v{r}ina\index{Barina, D.@Ba\v{r}ina, D.}'s verification to $2^{71}$~\cite{barina2025}, uses a sharper accounting of the
cycle structure than the inequality $\log(1+1/(3x_i))\le1/(3M)$.
It keeps the form of the cycle inequality: his bound~\cite[Theorem~27]{hercher2023} is $0<s/r-\lt<1/(4X\log2)$, in place of $1/(3X\log2)$, so the least length it does not exclude is again a rung, by the mediant lemma\index{mediant lemma} of \cref{chap:cycles} (Hercher finds it with his lemma on the least denominator of a fraction in an interval~\cite[Lemma~22]{hercher2023}), and every threshold of \cref{tab:ch8:thresholds} falls by the factor $\frac34$; his computer assisted Corollary~29 lowers the threshold for the rung $72\,057\,431\,991$ further, to $3\cdot2^{69}$.
Extend that computation to the later rungs: for each rung from $137\,528\,045\,312$ on, find the least verification range at which his accounting excludes it.
\end{problem*}

\section{Switching MM}\label{sec:ch11:smm}

\begin{problem*}[S1. Rates beyond quadratics]\designation{S1}\label{prob:ch11:s1}
Settled for the distance by \cref{thm:ch5b:rate}. For $\mu$ strongly convex and $L$ smooth $f$, with $\kappa=L/\mu$, the even iterates of $\smmdet$ contract the distance to the minimizer by $q=(1-1/\kappa)\max(1-(1+\beta)/\kappa,\beta)$ per period, which is the exact worst case over the class: explicit one dimensional functions attain it for every $\kappa$, and performance estimation\index{performance estimation} confirms it. The function gap contracts by $q^2$, which is exact for $\kappa\ge(1+\beta)/(1-\beta)$, where the quadratic $\frac\mu2\|x\|^2$ attains it. The remaining question is the function gap for $\kappa<(1+\beta)/(1-\beta)$, where performance estimation gives a smaller worst case, for example $0.049$ against $q^2=0.086$ at $\beta=\lt-1$ and $\kappa=2$: find it in closed form and describe the worst function.
\end{problem*}

\begin{problem*}[S2. Stochastic analysis]\designation{S2}\label{prob:ch11:s2}
Extend \cref{thm:smmvar} from quadratics to $\mu$ strongly convex, $L$ smooth objectives under
a bounded variance oracle: show that the stationary mean square error of
$\smmdet$ is at most $\tfrac53$ times that of gradient descent with the same
step, up to terms vanishing with the step size, for every $\beta<1$.
\end{problem*}

Partial progress. On quadratics the result holds for any uncorrelated noise,
and \cref{prop:ch7:mult} extends it to multiplicative (curvature) noise: under
the least squares scaling the slow mode ratio is unchanged, and under the
diffusive scaling a linear periodic method satisfying the hypotheses of \cref{thm:ch7:gainweight}, with gradient gains $c_i$ (\cref{def:ch7:gains}), has the mean square stability threshold $\epsilon^*=2\sum_ic_i/\sum_ic_i^2$, which is $1.472$ for $\smmdet$ against
$0.830$ for every step momentum\index{momentum} at $\beta=\lt-1$. What remains open is the
nonquadratic case and large multiplicative noise\index{multiplicative noise} at large step, where the
bound $5/3$ fails.

\begin{problem*}[S3. The best period two ratio]\designation{S3}\label{prob:ch11:s3}
Settled for quadratics by \cref{prop:ch7:pareto}. Every period two ratio $r$ between $1$ and the Chebyshev ratio $r_{\mathrm{ch}}(\kappa)$,
correctly scaled, lies on the rate and noise frontier
$R=1-4\kappa/((2-N)(\kappa+1)^2)$, with noise floor\index{noise floor} $N=2-1/G(r)$ and
$G(r)=(1+r)^2/(4r)$, and larger ratios are dominated by the Chebyshev ratio in both criteria; the Chebyshev\index{Chebyshev step sizes} ratio is the fast end, gradient descent the
quiet end, and the SMM ratio $1+\beta$ sits near the quiet end with a $5.4$
percent gain in the rate exponent for a $5.1$ percent rise in noise. The open part is the
choice of scale without knowledge of $\mu$: the optimal scale is
$\alpha L\to(2+\beta)/(1+\beta)$, not $1$.
\end{problem*}

\begin{problem*}[S4. Adaptive switching]\designation{S4}\label{prob:ch11:s4}
$\smmstate$ suppresses injection when progress is adequate. Find a principled
switching rule, using only $L$ and observed function values, that retains the
uniform noise bound of \cref{prop:ch7:bound} while closing the noiseless gap
to every step momentum identified in \cref{chap:noise}.
\end{problem*}

Partial progress. On quadratics, at every parameter value computed, no switching rule is unstable or much slower than the worse of gradient descent and every step Nesterov\index{Nesterov acceleration}: the joint spectral radius of the descent and injection matrices, computed exactly by \cref{prop:ch5b:jsr}, is below one, and at $\beta=\lt-1$ and $\beta=0.9$ it exceeds the larger pure rate by at most $0.006$.
A proof of the transient bound \cref{conj:ch5b:transient} would make the stability unconditional, at all parameter values. The noise side of the problem, a rule whose stationary variance is
provably bounded by $\frac53$ of gradient descent, remains open.

\section{Computational projects}\label{sec:ch11:projects}

The following projects are finite, well specified, and suited to a
collaborator with computing resources. Each comes with a clear success
criterion.

\paragraph{P1. A machine checked cycle bound.} Formalize the bound of
\cref{thm:cycles} in Lean\index{Lean}. The affine identity and the inequality $3^m<2^k$ are
already proved in Lean~4 using only its core library (\cref{sec:ch8:lean}), and
\cref{sec:ch8:certificate} gives an exact integer certificate\index{certificate}: an enclosure of
$\log_23$ by two convergents\index{convergent} certified with $320$ bit mantissas, an enclosure
of $\log2$, the $44$ ladder rungs up to $72\,057\,431\,991$, and an exclusion
inequality at each rung using integers of at most $248$ bits. What remains is
to formalize the upper half of the cycle inequality in the logarithm free form
$2^k(3X)^m\le3^m(3X+1)^m$, the mediant lemma\index{mediant lemma}, \cref{thm:ch8:onesided} and a verified squaring routine,
and to replay the certificate. Success means a Lean file, with no
\texttt{sorry}, whose main theorem states that every nontrivial cycle of the accelerated map on the positive integers has at least $72\,057\,431\,991$ odd steps, conditional only on a stated hypothesis encoding the verified range\index{verification range}: every positive integer below $2^{68}$ reaches $1$.

\paragraph{P2. Thresholds and the window.} The verification thresholds at which
the sharp bound moves are now computed for $60\le\log_2X\le90$
(\cref{tab:ch8:thresholds}); the next step past the current range $2^{71}$~\cite{barina2025} is at $2^{71.884}$. The
open computational question is the window of \cref{lem:ch8:window}: for the
hypothetical cycles at the rungs, compute the largest lower bound for the least
element that can be proved from the structure of $c_k(w)$ for words with the
prescribed $(m,k)$ that satisfy a further restriction, for instance compatibility with the first steps of an orbit starting above $2^{71}$.
Without a restriction nothing beyond the trivial bound $3^m-2^m$ for $c_k(w)$ is possible, since the word $1^m0^{k-m}$ attains it.

\paragraph{P3. Orbits from very large starting values.} The study of
\cref{sec:ch9:single} followed orbits from near $2^{10000}$ with residues
modulo $3^5$. Extend it to starting values near $2^{10^5}$ and beyond, to moduli
$3^a$ with $a\le8$, to joint statistics of consecutive residues, and to the
dyadic digits of the iterates, and look for any systematic departure from
the model chain of \cref{lem:ch9:mixing}.

\paragraph{P4. Total stopping time cumulants at scale.} Extend the exact census of
\cref{sec:ch3:missing} beyond $2^{30}$ (memory permitting) and the sampled
version to $10^{15}$, test whether the additive variance deficit saturates,
and fit the third cumulant.

\paragraph{P5. Switching MM at scale.} Run $\smmdet$ against every step
Nesterov\index{Nesterov acceleration} and heavy ball\index{heavy ball} at $\beta\in\{0.585,0.9\}$, and against Adam, on
minibatch training of standard models where the gradient noise is genuine,
with the base step size of gradient descent and the three momentum methods fixed at the common value $\alpha=1/L$, at which all four are stable and every eigendirection of the Hessian has $a=\alpha\lambda\le1$.
Report final loss and its variance over seeds, and compare the measured noise floors with the prediction of \cref{thm:smmvar} evaluated on the spectrum of the Hessian at the solution, whose ratio to gradient descent lies between $1$ and $(2+2\beta+\beta^2)/(2+\beta)$ on every eigendirection (\cref{prop:ch7:uniform}). Larger common steps are possible only below $2(1+\beta)/((1+2\beta)L)$, the stability limit of every step Nesterov, where the denominator of \eqref{eq:ch7:nagvar} vanishes and which lies below the limit $(2+\beta)/((1+\beta)L)$ of $\smmdet$ (\cref{cor:ch6:stability}); above $a=1$ \cref{prop:ch7:uniform} does not apply; \eqref{eq:ch7:vavg} still holds there, since the proof of \cref{thm:smmvar} uses only $|\rho|<1$, and at $\beta=0.585$ it gives a ratio above $(2+2\beta+\beta^2)/(2+\beta)$ once $a>1.41$, reaching $1.55$ at the Nesterov limit.

\paragraph{P6. Periodic schedules in general.} Settled in the slow mode by the
gain weight\index{gain weight theorem} theorem (\cref{thm:ch7:gainweight}): no stable linear periodic
first order method beats gradient descent at matched speed, and the penalty is
the Cauchy and Schwarz ratio $p\sum c_i^2/(\sum c_i)^2$ of the gradient gains.
What remains computational is the regime away from the slow mode: for a given
spectrum and noise level, find the period $p$ schedule minimizing the summed
stationary variance at a prescribed worst case rate, and compare it with the
period two frontier of \cref{prop:ch7:pareto}.

\paragraph{P7. Erd\H{o}s $125$ beyond $4^{13}$.} Extend the exact density
computation of \cref{chap:cycles} to $4^{15}$ or beyond with a bit array of a
few gigabytes, and test whether the successive minima just below powers of $3$
track the quality of the approximations $3^m\approx4^k$ quantitatively.

\paragraph{P8. Perfect samplers beyond the coefficient stopping time.} The sampler of
\cref{thm:ch10b:perfect} is exact because the level sets of the coefficient stopping time $\tau$ are unions
of residue classes. Find other energies with this property whose level sets admit a dynamic programme of size polynomial in the depth,
for instance the number of descents in the first $k$ steps or the coefficient stopping time
of $T^{(j)}$ for fixed $j$, give the dynamic programme for each, and use the
samplers to measure the frequency with which the block potentials of
\cref{prop:ch5c:runpot} fail at depth $k$, as a function of $k$.

\paragraph{P9. Trees of slow integers.} \Cref{tab:ch10b:clusters} shows that
the $257$ integers below $2^{26}$ with total stopping time above $400$, counted in steps of $T$ as in \cref{chap:samplers}, are $34$
trees of merged orbits, and the $24$ above $440$ are six. \Cref{tab:ch10b:trees} counts the trees
of $\{\sigma_\infty>m\}\cap[2,2^B)$ for $B=22$, $24$ and $26$: over this range their number grows with $B$
roughly like $2^{cB}$ with $c$ between $1$ and $1.4$, and their mean size like
$2^{0.4B}$.
These exponents are local: at fixed $m$ each tree has its own root, an integer with $\sigma_\infty=m-100$, and there are only finitely many such integers, all at most $2^{m-100}$ since a step of $T$ at most halves, so the number of trees is bounded in $B$.
Explain the exponents observed for $22\le B\le26$, or find how the number and the size of the trees grow when $m$ and $B$ increase together, for instance with $m/B$ fixed.
The number of trees is the
number of points at remaining depth $m-100$ reached from the box by an orbit segment of more than a hundred steps, a statistic of the merging mechanism of \cref{lem:ch10b:merge} that has not, to our knowledge, been studied; the mean size is a count on the inverse tree of \cref{chap:difference}, and whether the certified exponents of that chapter bear on it, when, at the deeper levels, the roots usually lie above the box, is part of the problem. Relate the
roots of the trees to the delay records of \cref{chap:intro}.

\section*{Notes}

Problems O1, O3, C2 and the cycle formulation are restatements of questions
going back to Terras\index{Terras, R.}~\cite{terras1976}, Steiner\index{Steiner, R. P.}~\cite{steiner1977} and
Eliahou\index{Eliahou, S.}~\cite{eliahou1993}, sharpened by the results of this book. Problem C1
and the Erd\H{o}s $125$ analogy come from our note on the AlphaProof Nexus\index{AlphaProof Nexus}
system~\cite{polsonsokolov2026nexus,tsoukalas2026}. Problems S1 to S4 and P5,
P6 grew out of the analysis of \cref{chap:smm,chap:noise}, and P8 and P9 out
of \cref{chap:samplers}.
